\chapter{Assignment 1 --- Simulator Installation}
\label{ch:hw1}

\section{Environment}
\label{sec:hw1-env}

All work was done on a MacBook Pro with an Apple M3 Max processor
(\texttt{arm64}), running macOS with Homebrew~6.0.21 and CMake~4.4.3. The CMake
version matters: it is the source of most of the build failures recorded below.

\section{Webots}
\label{sec:hw1-webots}

Webots is distributed both as a disk image and as a Homebrew cask
(\url{https://cyberbotics.com/}). Either path works; the cask was used.

\begin{lstlisting}[style=shell,caption={Installing Webots.}]
# option 1 - download the disk image and run it
# https://github.com/cyberbotics/webots/releases/download/R2025a/webots-R2025a.dmg

# option 2 - Homebrew cask
brew install --cask webots
\end{lstlisting}

\begin{figure}[ht]
\centering
\includegraphics[width=0.85\linewidth]{hw1_webots_install.png}
\caption{Webots installed on macOS.}
\label{fig:hw1-webots-install}
\end{figure}

\begin{figure}[ht]
\centering
\includegraphics[width=0.85\linewidth]{hw1_webots_running.png}
\caption{Webots running a sample world.}
\label{fig:hw1-webots-running}
\end{figure}

\section{ARGoS}
\label{sec:hw1-argos}

ARGoS is documented at \url{https://www.argos-sim.info/} and is built from
source at \url{https://github.com/ilpincy/argos3}.

\subsection{The Homebrew tap fails}
\label{sec:hw1-argos-brew}

The maintainer publishes a tap, and the tap owner and the repository owner are
the same account, so the tap looked like the safe route.

\begin{lstlisting}[style=shell,caption={The tap route.}]
brew tap ilpincy/argos3
brew trust ilpincy/argos3
brew install bash-completion qt argos3
\end{lstlisting}

It fails during CMake configuration:

\begin{lstlisting}[style=shell,caption={CMake rejects the project's minimum version.}]
CMake Error at CMakeLists.txt:15 (cmake_policy):
  Compatibility with CMake < 3.5 has been removed from CMake.

  Update the VERSION argument <min> value.  Or, use the <min>...<max> syntax
  to tell CMake that the project requires at least <min> but has been updated
  to work with policies introduced by <max> or earlier.

  Or, add -DCMAKE_POLICY_VERSION_MINIMUM=3.5 to try configuring anyway.

-- Configuring incomplete, errors occurred!
\end{lstlisting}

The cause is the CMake version: the local CMake is~4.4.3, and ARGoS still
declares compatibility with CMake~$<3.5$, which CMake~4 no longer accepts. The
error message names its own fix --- \texttt{-DCMAKE\_POLICY\_VERSION\_MINIMUM=3.5}
--- but the tap gives no place to pass it. So the tap was removed and ARGoS was
built by hand instead.

\begin{lstlisting}[style=shell,caption={Removing the tap.}]
cmake --version      # cmake version 4.4.3
brew untap ilpincy/argos3
\end{lstlisting}

\subsection{Building from source}
\label{sec:hw1-argos-source}

Building by hand makes it possible to pass the policy flag, and
\path|-DCMAKE_INSTALL_PREFIX=$HOME/.local| keeps the install in user space
so no \texttt{sudo} is needed.

\begin{lstlisting}[style=shell,caption={Dependencies and source build.}]
brew install pkg-config cmake libpng freeimage lua qt docbook asciidoc \
             graphviz doxygen

git clone https://github.com/ilpincy/argos3.git argos3
cd argos3
mkdir build_simulator
cd build_simulator

cmake ../src -DCMAKE_POLICY_VERSION_MINIMUM=3.5 \
             -DCMAKE_INSTALL_PREFIX=$HOME/.local \
             -DCMAKE_PREFIX_PATH=$(brew --prefix qt)
make
make doc
make install
\end{lstlisting}

Two failures were hit on the way to that command, and both are already fixed in
it:

\begin{enumerate}
  \item \textbf{\texttt{make install} cannot find \texttt{argos3.1.gz}.} The man
        page is a build target of its own. Running \texttt{make doc} before
        \texttt{make install} produces it.
  \item \textbf{\texttt{find\_package(ARGoS)} cannot find
        \texttt{FindARGoSQTOpenGL.cmake}} when building the examples. The Qt
        OpenGL visualization was not built, because CMake did not find Qt.
        Adding \texttt{-DCMAKE\_PREFIX\_PATH=\$(brew --prefix qt)} and
        rebuilding fixes it.
\end{enumerate}

\subsection{Verifying the installation}
\label{sec:hw1-argos-verify}

Because the prefix is \texttt{\$HOME/.local}, its \texttt{bin} directory must be
on \texttt{PATH}; otherwise \texttt{argos3} reports \texttt{command not found}
even though it installed correctly.

\begin{lstlisting}[style=shell,caption={Putting the binary on PATH.}]
echo 'export PATH="$HOME/.local/bin:$PATH"' >> ~/.zshrc
source ~/.zshrc

argos3 -v
\end{lstlisting}

\subsection{Running an example}
\label{sec:hw1-argos-example}

The example worlds are listed at \url{https://www.argos-sim.info/examples.php}.

\begin{lstlisting}[style=shell,caption={Building the examples.}]
cd argos3-examples
mkdir build
cd build
cmake ../ -DCMAKE_POLICY_VERSION_MINIMUM=3.5 \
          -DCMAKE_PREFIX_PATH=$HOME/.local
make
\end{lstlisting}

\begin{figure}[ht]
\centering
\includegraphics[width=0.85\linewidth]{hw1_argos3_running.png}
\caption{ARGoS running an example experiment.}
\label{fig:hw1-argos-running}
\end{figure}

\subsection{Uninstalling}
\label{sec:hw1-argos-uninstall}

The install writes its own uninstall script --- the install log lists
\path|$HOME/.local/share/argos3/uninstall_argos3.sh|. That is the official
removal path and it should be preferred.

\begin{lstlisting}[style=shell,caption={Uninstalling ARGoS.}]
bash ~/.local/share/argos3/uninstall_argos3.sh

# if the script is missing, remove the installed paths by hand.
# no sudo needed - everything is under the ~/.local prefix.
rm -f  ~/.local/bin/argos3
rm -rf ~/.local/include/argos3
rm -rf ~/.local/lib/argos3
rm -f  ~/.local/lib/libargos3*
rm -rf ~/.local/share/argos3

which argos3        # no output means the removal is complete
\end{lstlisting}
